% Generated draft: requires mathematical and source review.
\documentclass[11pt]{article}
\usepackage[margin=25mm]{geometry}
\usepackage{fontspec}
\setmainfont{Latin Modern Roman}
\usepackage{amsmath,amssymb,mathtools,unicode-math}
\setmathfont{Latin Modern Math}
\usepackage{enumitem,needspace,xurl,hyperref}
\setlength{\parindent}{0pt}
\setlength{\parskip}{5pt}
\setlength{\emergencystretch}{3em}
\newcommand{\N}{\mathbb{N}}
\newcommand{\Z}{\mathbb{Z}}
\newcommand{\Q}{\mathbb{Q}}
\newcommand{\R}{\mathbb{R}}
\newcommand{\C}{\mathbb{C}}
\newcommand{\F}{\mathbb{F}}
\newcommand{\A}{\mathbb{A}}
\newcommand{\PP}{\mathbb{P}}
\newcommand{\E}{\mathbb{E}}
\newcommand{\eps}{\varepsilon}
\newcommand{\dd}{\,\mathrm{d}}
\begin{document}
\section*{UnsolvedMath ID: 10400084}
\textbf{Source problem code:} catalogue entry
\subsection{Definitions and mathematical statement}
Problem 4.6 asks to calculate ``the skein homology based on the Kauffman bracket skein relation.'' In a usual bracket convention the relation is
\[
 [L_{\times}]=A[L_0]+A^{-1}[L_{\infty}],
\]
with the three local diagrams differing by a crossing and its two smoothings; the usual trivial-circle relation is \([L\sqcup O]=-(A^2+A^{-2})[L]\). These formulas explain the bracket terminology, but they define skein relations rather than a unique chain complex.

To identify a homology problem one must specify chain modules \(C_i\), differentials \(d_i:C_i\to C_{i-1}\) with \(d_i d_{i+1}=0\), and the intended input objects. The requested groups would then be
\[
 H_i(C_\bullet)=\ker(d_i)/\operatorname{im}(d_{i+1}),
\]
with any additional gradings or module structures. The supplied record contains no chain construction, differential, coefficient ring, bracket parameter, manifold or link class, degree range, or output criterion. Its definite article ``the'' points to missing surrounding definitions. Consequently the precise skein homology to be calculated remains unidentified. The research summary's mention of a relation to Khovanov homology does not specify that this is Khovanov homology or supply an identification with it.
\subsection{Short English statement}
Compute the homology theory intended by the source's Kauffman-bracket skein construction. The skein relation alone does not determine that theory, and the defining complex is absent. The supplied open-triage assessment only mentions a relation with Khovanov homology.
\subsection{Sources}
\begin{itemize}
\item Local UnsolvedMath registry, permanent ID \(10400084\), exact source code \texttt{catalogue entry}, Problem 4.6. No individual proposer or proposed year is supplied.
\item Ohtsuki, \emph{Problems on invariants of knots and 3-manifolds} (2002), Problem 4.6, PDF page 75: \url{https://msp.org/gtm/2002/04/gtm-2002-04-024s.pdf}. Bold printed headings distinguish 231 numbered items from surrounding references.
\item Extraction pdf-font-structure, accessed 2026-07-29. Source evidence NEEDS\_REVIEW presents an open problem and requires a current primary-literature audit. Supplied research classification OPEN-TRIAGE says the question concerns Kauffman-bracket skein homology and is related to Khovanov homology, without a further citation or definition.
\item Dataset attribution: Ulam AI / UnsolvedMath Contributors, CC BY 4.0. Source rights NEEDS\_REVIEW: publicly accessible author/publisher copy, redistribution terms requiring review.
\end{itemize}
\end{document}
